\input{build/foundations-preamble.tex}
\usepackage[authoryear]{natbib}
\renewcommand{\refname}{સંદર્ભગ્રંથો}
\renewcommand{\partname}{ભાગ}
\newenvironment{intro}{}{}
\newcommand{\Id}[1]{\mathord{\mathrm{Id}_{#1}}}
\newcommand{\equivrep}[2]{[#1]_{#2}}
\newcommand{\equivclass}[2]{#1/_{\!{#2}}}
\newcommand{\funrestrictionto}[2]{#1\mathord{\restriction}_{#2}}
\newcommand{\funimage}[2]{#1[#2]}
\newcommand{\emptyseq}{\Lambda}
\newcommand{\liff}{\mathbin{\leftrightarrow}}
\hypersetup{pdftitle={ગણો અને સંબંધો — ઓપન લોજિક ગુજરાતી}}
\begin{document}
\begin{titlepage}
{\Huge\bfseries ગણો અને સંબંધો\par}
\vspace{1em}
{\Large ઓપન લોજિક — ગુજરાતી\par}
\vspace{2em}
ઓપન લોજિક પ્રોજેક્ટના અંગ્રેજી મૂળનો ગુજરાતી અનુવાદ.\par
\vspace{1em}
આ આવૃત્તિમાં ગણો અને સંબંધોનાં સંપૂર્ણ પ્રકરણો છે.
ગણો, ઉપગણો, કાર્તેઝીય ગુણાકાર અને રસેલના વિરોધાભાસ પછી,
સંબંધોના ગુણધર્મો, સામ્ય વર્ગો, ક્રમો, આલેખો, વૃક્ષો અને
સંબંધો પરની ક્રિયાઓનો અભ્યાસ થાય છે.\par
\vfill
આ યંત્ર દ્વારા કરેલો અનુવાદ છે. સંપૂર્ણ ગ્રંથનું કામ ચાલુ છે.
આ આવૃત્તિમાં મૂળનાં ૧૬ એકમોનો સમાવેશ છે; કુલ ૭૨૨ એકમો છે.
કેટલીક વિશિષ્ટ પરિભાષા કામચલાઉ છે. ચકાસણીની વિગતો અને
મૂળ પાઠ વિશેની નોંધો સંપાદકીય માહિતીમાં આપેલી છે.\par
\vspace{1em}
\href{https://github.com/OpenLogicProject/OpenLogic}{Open Logic Project}
\quad \href{https://creativecommons.org/licenses/by/4.0/}{CC BY 4.0}\par
\href{https://github.com/KokunoYumeto/OpenLogic-translations}{અનુવાદોનું કેન્દ્ર}\par
\end{titlepage}
\tableofcontents
\clearpage
\input{build/foundations-body.tex}
\clearpage
\section*{સંપાદકીય નોંધો}
\addcontentsline{toc}{section}{સંપાદકીય નોંધો}
આ નોંધો મૂળ પાઠથી અલગ છે. મૂળની ગણિતીય નિશાનીઓ જાળવી છે.
સંબંધોને ગણ તરીકે રજૂ કરતા ઉદાહરણમાં વિકર્ણ માટે \(I\) વપરાયું છે,
પણ તેને અલગથી વ્યાખ્યાયિત કરેલું નથી. વૃક્ષની શાખાની વ્યાખ્યામાં
\(X\) વપરાયું છે, જ્યારે વૃક્ષનો આધારગણ \(A\) છે; અહીં મૂળનો
\(X\) જાળવ્યો છે. ઉપવૃક્ષના ઉદાહરણમાં ખાલી ઉપગણનો અપવાદ
સ્પષ્ટ નથી, જોકે મૂળવાળા વૃક્ષની વ્યાખ્યા તેને બાકાત રાખે છે.
ક્રમોના વિભાગમાં \(R^+\) સ્વવાચક સંવરણ દર્શાવે છે અને
સંબંધોની ક્રિયાઓના વિભાગમાં તે પરંપરિત સંવરણ દર્શાવે છે;
દરેક વિભાગની સ્થાનિક વ્યાખ્યા લાગુ પડે છે.

આવૃત્તિમાં ગણો અને સંબંધોનાં બંને પ્રકરણો સાથે વાંચી શકાય છે.
પછીનાં પ્રકરણોના સશરતી ઉલ્લેખોમાં મૂળે આપેલો વૈકલ્પિક પાઠ
વપરાયો છે. પરિભાષા, મૂળ સ્રોતો અને ચકાસણીની વધુ વિગતો
આવૃત્તિની સાથેની ફાઇલોમાં છે.
\begin{thebibliography}{1}
\bibitem[Benacerraf(1965)]{Benacerraf1965}
Benacerraf, Paul. 1965.
\emph{What numbers could not be}.
The Philosophical Review 74(1), 47--73.
\end{thebibliography}
\end{document}
